\subsection*{1.4. 变差}

\subsubsection*{1.4.1.}
设 $S$ 为 0.2.5 中的 Hensel 离散赋值环谱, $Y$ 为 $s$ 上的概形. 设 $\Lambda$ 为环. 更一般地, 可取 $\Lambda$ 为 $Y$ 上的环层. 对 $Y\times_sS$ 上的 $\Lambda$-模层所成范畴的导出范畴 $D(Y\times_sS,\Lambda)$ 的每个对象 $K$, 我们要定义一个\underline{变差}态射.

\subsubsection*{1.4.2.}
$Y\times_sS$ 上的 $\Lambda$-模层复形 $K$ 等同于如下三元组 $(K_s,K_\eta,\varphi)$:
\begin{enumerate}[label=\alph*)]
\item $K_s$（分别 $K_\eta$）是 $\bar Y$ 上的层复形 $K_{\bar s}$（分别 $K_{\bar\eta}$）, 带有 $\operatorname{Gal}(\bar s/s)$（分别 $\operatorname{Gal}(\bar\eta/\eta)$）的连续作用;

\item $\varphi$ 为等变态射 $\varphi:K_{\bar s}\longrightarrow K_{\bar\eta}$.
\end{enumerate}
这里及以下均使用 1.4.1 中固定的概形 $Y$; 原印本 1.4.2 的首句误写为 $X\times_sS$.

这样的复形 $K$ 总与某个复形 $K'=(K'_s,K'_\eta,\varphi')$ 同伦等价, 其中 $\varphi'$ 单射, 且正合列
\begin{equation}\tag{1.4.2.1}
0\longrightarrow K'_{\bar s}
\longrightarrow K'_{\bar\eta}
\longrightarrow\operatorname{coker}(\varphi')
\longrightarrow0
\end{equation}
逐次分裂. 只须取 $K'_s=K_s$, 并取 $K'_\eta$ 为 $K_\eta$ 与 $\mathrm{sp}^*K_s$ 的锥的直和. 置
\begin{equation*}
\Phi(K)=\operatorname{coker}(\varphi').
\end{equation*}
正合列 (1.4.2.1) 在 $Y\times_s\eta$ 上的 $\Lambda$-模层复形的同伦范畴 $K(Y\times_s\eta,\Lambda)$ 中定义一个区别三角
\begin{equation}\tag{1.4.2.2}
\begin{tikzcd}[column sep=large, row sep=large]
{} & \Phi(K) \arrow[dl] & {} \\
\mathrm{sp}^*K_s \arrow[rr] & {} & K_\eta \arrow[ul]
\end{tikzcd}
\end{equation}
此三角只依赖于 $K\in\operatorname{Ob}K(Y\times_sS,\Lambda)$.

这个构造可下降到导出范畴, 对 $D(Y\times_sS,\Lambda)$ 的每个对象 $K$ 给出 $D(Y\times_s\eta,\Lambda)$ 中的一个区别三角.

三角 (1.4.2.2) 主要通过它所诱导的长正合超上同调列发挥作用:
\begin{equation*}
\cdots\longrightarrow H^i(\bar Y,K_{\bar s})
\longrightarrow H^i(\bar Y,K_{\bar\eta})
\longrightarrow H^i(\bar Y,\Phi(K)_{\bar\eta})
\longrightarrow\cdots.
\end{equation*}

\subsubsection*{1.4.3.}
惯性群 $I$ 在 $K'_{\bar s}$ 上平凡作用. 因此, 对每个 $\sigma\in I$, $K'_{\bar\eta}$ 的自同态 $\sigma-1$（记号见 (1.4.2.1)）通过一个态射
\begin{equation*}
\operatorname{Var}(\sigma):
\operatorname{coker}(\varphi')\longrightarrow K'_{\bar\eta}.
\end{equation*}
分解. 这个构造可下降到导出范畴: 对 $K\in\operatorname{Ob}D(Y\times_sS,\Lambda)$ 及 $\sigma\in I$, 它定义\underline{变差}
\begin{equation}\tag{1.4.3.1}
\operatorname{Var}(\sigma):
\Phi(K)_{\bar\eta}\longrightarrow K_{\bar\eta}.
\end{equation}
以 $q$ 表示从 $K_\eta$ 到 $\Phi(K)$ 的箭头, 有
\begin{align}\tag{1.4.3.2}
\sigma&=1+\operatorname{Var}(\sigma)q
&&\text{作用于 }K_{\bar\eta},\\
\tag{1.4.3.3}
\sigma&=1+q\operatorname{Var}(\sigma)
&&\text{作用于 }\Phi(K)_{\bar\eta},
\end{align}
及
\begin{equation}\tag{1.4.3.4}
\operatorname{Var}(\sigma\tau)
=\operatorname{Var}(\sigma)q\operatorname{Var}(\tau)
 +\operatorname{Var}(\sigma)+\operatorname{Var}(\tau).
\end{equation}
